%%%%%%%%%%%%%%%%%%%%%%%%%%%%%%%%%%%%%%%%%%%%%%%%%%%%%%%%%%%%%%%%%%%%%%%%
% Plantilla TFG/TFM
% Escuela Politécnica Superior de la Universidad de Alicante
% Realizado por: Jose Manuel Requena Plens
% Contacto: info@jmrplens.com / Telegram:@jmrplens
%%%%%%%%%%%%%%%%%%%%%%%%%%%%%%%%%%%%%%%%%%%%%%%%%%%%%%%%%%%%%%%%%%%%%%%%

\chapter*{Preámbulo}
\thispagestyle{empty}
\vspace{1cm}
Este proyecto consiste en el desarrollo de un sistema de Inteligencia Artificial generativa para lectura facilitada, que tiene como objetivo facilitar 
la comprensión de textos a personas con dificultades lectoras. Este sistema utiliza distintas técnicas avanzadas de procesamiento del lenguaje natural y aprendizaje 
automático para generar resúmenes y explicaciones adaptadas a las necesidades de usuarios con diferentes niveles de comprensión. El sistema se basa en modelos de lenguaje 
ya entrenados con una amplia variedad de textos y utiliza algoritmos de generación de texto para crear resúmenes y explicaciones. El objetivo final es proporcionar 
una herramienta útil para mejorar la comprensión lectora y facilitar el acceso a la información.


\cleardoublepage %salta a nueva página impar
\chapter*{Agradecimientos}

\thispagestyle{empty}
\vspace{1cm}

En primer lugar, este trabajo no habría sido siquiera posible sin la ayuda y dirección de mi tutor, José Abreu Salas, bajo cuya supervisión escogí y desarrollé el tema a estudiar. Él ha sido quien me ha guiado durante toda la realización de este trabajo.
\\
\par También quiero agradecer a mi tutor Rafael Muñoz Guillena, y en general a todo el Grupo de Procesamiento del Lenguaje Natural y Sistemas de Información (GPLSI), por haberme permitido realizar las prácticas externas en sus instalaciones, por su apoyo y por la ayuda que me han brindado en la elaboración de este trabajo.
\\
\par Finalmente, quiero mostrar mi gratitud a mi familia y a los amigos que he conocido durante la carrera, quienes han sido mi mayor apoyo a lo largo de estos años. Sin ellos no habría sido capaz de llegar hasta aquí.


\cleardoublepage %salta a nueva página impar
% Aquí va la cita célebre si la hubiese. Si no, comentar la(s) linea(s) siguientes
\chapter*{}
\setlength{\leftmargin}{0.5\textwidth}
\setlength{\parsep}{0cm}
\addtolength{\topsep}{0.5cm}
\begin{flushright}
\small\em{
``No todos los que vagan están perdidos"
}
\end{flushright}
\begin{flushright}
\small{
J.R.R. Tolkien.
}
\end{flushright}
\cleardoublepage %salta a nueva página impar
